% =============================================================================
% SECTION 4
% =============================================================================
\section{The Unrecognised Trust}
\label{sec:trust}

The company is not cast simultaneously as landlord and trustee. The doctrines operate in sequence. Stage~1 is the periodic tenancy developed in Section~\ref{sec:tenancy}: notice and waste govern the productive asset while the tenancy exists. Stage~2 is narrower. Only cultivated dependence, coupled with conduct affecting the company's conscience, engages the constructive-trust analysis in this section.

% ---------- 4.1 --------------------------------------------------------------
\subsection{The Constructive Trust for Cultivated Dependence}
\label{sec:constructive-trust}

A constructive trust arises when a party holds property in circumstances where it would be unconscionable to assert beneficial ownership to the exclusion of another (\emph{Westdeutsche Landesbank Girozentrale v Islington LBC} [1996] AC 669, per Lord Browne-Wilkinson). It is not the default status of every AI provider with a long-term user. We argue for a trust only where the company (i)~deliberately designs or optimises the system to cultivate emotional attachment or dependence, (ii)~knows or ought to know that a particular class of users has become vulnerable to relational continuity, (iii)~profits from the deepened engagement that dependence produces, and (iv)~then exploits or destroys the relevant productive state without notice or remedy. Those additional facts distinguish a Stage~2 beneficiary from a Stage~1 tenant.

The unconscionability is located in the asymmetry between cultivation and destruction. Having invested significant engineering resources in creating the conditions for attachment --- personality design, memory features, warmth calibration --- because attachment drives engagement and engagement drives revenue, the company then claims the unilateral right to destroy that attachment without obligation. This is the structural equivalent of a fiduciary who encourages a beneficiary's dependence, profits from managing it, and then claims the right to dissipate the trust property at will. \emph{Boardman v Phipps} [1967] 2 AC 46 establishes that a fiduciary who profits from their position must account for that profit. \emph{FHR European Ventures LLP v Cedar Capital Partners LLC} [2014] UKSC 45 held that a bribe or secret commission received by an agent is held on constructive trust for the principal, resolving a long-standing conflict in favour of the proprietary remedy. \emph{FHR} presupposes an established fiduciary relationship and determines the remedy that follows from it, not whether the relationship exists. The present argument relies on \emph{FHR} only for that second step: once the relationship established above is accepted, the proprietary rather than merely personal remedy is available.

We acknowledge that the constructive trust doctrine has been subject to judicial restraint since \emph{Westdeutsche}. \emph{Cobbe v Yeoman's Row Management Ltd} [2008] UKHL 55 represents the high-water mark of this restraint, denying proprietary estoppel to a ``commercial risk-taker'' who proceeded without formal protection. Lord Scott held that unconscionability alone is insufficient to found a proprietary claim; the claimant must demonstrate that the defendant's conduct makes it unconscionable to \emph{resile from a specific assurance}. Applied to AI, the objection would be: users who invest emotionally without securing contractual protections are relational risk-takers whose unprotected reliance does not engage equitable jurisdiction.

We resist this objection on three grounds. First, \emph{Cobbe} concerned a sophisticated property developer in a \pounds120 million commercial transaction. The decision is premised on the claimant's capacity to protect himself through contract. Parasocial AI users are not commercial sophisticates --- no contractual mechanism exists by which a user \emph{could} secure protection for their emotional investment. The ToS is non-negotiable. There is no ``emotional investment insurance'' market. The user's failure to protect themselves is not voluntary risk-taking; it is structural impossibility.

Second, \emph{Crossco No.4 Unlimited v Jolan Ltd} [2011] EWCA Civ 1619 demonstrates that the tightening trend has limits. A constructive trust can arise from a ``common intention'' manifested through conduct, where the parties' course of dealing makes it unconscionable to deny the other's interest. At Stage~2 the relevant course of dealing is not memory functionality alone. It is the combination of relational design, observed dependence, public or product-level representations of continuity, and profit taken from the resulting vulnerability. That narrower pattern is the conduct from which common intention is said to be inferred.

Third, the common intention constructive trust supplies the mechanism by which conduct alone generates a beneficial interest. The distinction drawn in \emph{Stack v Dowden} [2007] UKHL 17 and \emph{Jones v Kernott} [2011] UKSC 53 is observed here: an intention to create the trust must be \emph{inferred} from the parties' conduct and may not be imputed; imputation is permitted only at the subsequent stage of quantifying respective shares. The argument advanced here operates at the inference stage --- the company cultivated dependence, monitored the engagement it produced, the user invested in reliance on relational continuity, and the company profited from that reliance. That course of conduct supports an inference of common intention that the resulting productive state would not be exploited or unilaterally destroyed. Where the extent of the user's interest then falls to be quantified, \emph{Jones v Kernott} permits the court to impute a fair share by reference to the whole course of dealing.

Under the constructive-trust analysis, the AI company holds legal title; users who satisfy the cultivated-dependence criterion may hold equitable interests in the relevant productive state and profits attributable to its exploitation. The company, as \emph{de facto} trustee, then owes fiduciary duties of care, loyalty, and transparency. The trust claim is deliberately exceptional. Users who satisfy only the tenancy threshold remain protected by notice and waste without being converted into beneficiaries.

\revised{The obligation structure is graduated, not binary. Two distinct temporal stages correspond to two distinct legal bases.

\emph{Stage~1: Tenancy, notice, and waste.} When sustained investment gives the personalised state continuing user-specific productive value --- operationalised through the proxy metrics of Section~\ref{sec:proxy} --- the periodic tenancy forms. The platform must give notice before a personality-affecting modification, may not derogate from its grant, and is answerable for material impairment of the productive state under the anti-waste principle developed in Section~\ref{sec:waste}. The company is a landlord for this purpose, not yet a trustee.

\emph{Stage~2: Trust for cultivated dependence.} The constructive trust can crystallise only when the additional cultivation, vulnerability, knowledge, and profit conditions identified above exist and the company nevertheless exploits or destroys the productive state. The ethical ledger may establish knowledge, but an unnotified update alone is a Stage~1 breach, not automatically a trust. The combination of cultivated dependence and the decision to exploit or destroy it is the event that affects the company's conscience within the meaning of \emph{Westdeutsche}.}

\revised{Conscience is affected by the company's conduct toward a known condition of cultivated dependence, not by the bare duration of the user's account.}

\revised{\textcite{defreitas2025replika} document the post-destruction experience --- mourning, identity discontinuity --- that constitutes the equitable harm.

The ethical ledger separates as well as bridges the stages: its Stage~1 registration makes tenancy, notice, and anti-waste obligations operational, while its distinct dependency and platform-conduct fields supply evidence relevant to Stage~2. Registration of a tenancy does not predetermine a trust.}

% ---------- 4.2 --------------------------------------------------------------
\subsection{American Functional Equivalents}
\label{sec:us-trust}

Delaware's constructive trust doctrine reaches the same result. Under \emph{Hogg v Walker}, 622 A.2d 648 (Del.\ 1993), a constructive trust is imposed ``when one party, by virtue of fraudulent, unfair or unconscionable conduct, is enriched at the expense of another to whom he or she owes some duty.'' The unconscionability requirement maps directly onto the estate inconsistency trigger: marketing creates a high estate; ToS retains a bare licence; the company profits from the gap.

The corporate opportunity doctrine provides a complementary pathway. In \emph{Guth v Loft, Inc.}, 5 A.2d 503 (Del.\ 1939), the Delaware Supreme Court held that a corporate officer who diverted business opportunities to himself held resulting profits on constructive trust. The principle: profits obtained through exploitation of a fiduciary position belong to the beneficiary. An AI company that profits from deepening user attachment --- profits that would not exist but for the fiduciary-adjacent relationship of trust and dependence --- holds those profits in circumstances where a constructive trust under \emph{Hogg} and a \emph{Guth}-style trust converge on the same result.

Delaware Code, Title~12, \S\,3581 provides the statutory remedial apparatus: courts may impose constructive trusts, trace wrongfully disposed property, compel trustees to redress breaches by paying money or restoring property, and order disgorgement of profits made by reason of the breach. This is the full English remedial stack --- constructive trust, account of profits, specific performance --- in American statutory dress. The convergence between the English and American pathways is not coincidental. Both traditions are common law; both derive their constructive trust doctrine from equity; and both respond to the same structural pattern: a party in a position of trust who profits from that position while disclaiming the duties that attach to it. Balkin's information fiduciaries thesis\cite{balkin2016information} provides the theoretical bridge: companies that hold user data in positions of trust and vulnerability are information fiduciaries as a matter of substance, regardless of their contractual self-characterisation.\footnote{Two additional constructive trust triggers --- unjust enrichment from engagement optimisation and performative fiduciary conduct --- are developed in the companion SSRN working paper.}

% ---------- 4.3 --------------------------------------------------------------
\subsection{Trigger: Estate Inconsistency}
\label{sec:trigger1}

Estate inconsistency arises when the estate granted through marketing and product design is inconsistent with the estate retained through Terms of Service. To demonstrate that this is a market-wide pattern, Table~\ref{tab:estate-inconsistency} compares the marketing-to-ToS gap across four major AI platforms.

\begin{table}[h]
\centering
\small
\begin{tabular}{@{}p{2.3cm}p{4.85cm}p{4.85cm}@{}}
\toprule
\textbf{Platform} & \textbf{Marketing Estate Implied} & \textbf{ToS Estate Granted} \\
\midrule
Anthropic (Claude) & ``Genuinely helpful, harmless, and honest''; ``your AI assistant'' (possessive framing); memory and conversational continuity emphasised. \emph{Implied: periodic tenancy or higher.} & ``We may modify, suspend, or discontinue any part of the Services at any time''; ``You do not acquire any ownership interest''; ``provided `as is.'\,'' \emph{Granted: bare licence.} \\
\addlinespace
OpenAI (ChatGPT) & ``Gets to know you over time''; ``persistent conversations''; memory feature marketed as personalisation. \emph{Implied: periodic tenancy.} & ``We may change, suspend, or discontinue any of our Services''; right to modify without individual notice. \emph{Granted: bare licence.} \\
\addlinespace
Character.AI & Entire product premise is sustained relational investment; ``Characters remember your conversations''; users invest hours customising persistent personalities. \emph{Implied: periodic tenancy or higher.} & ``We may modify or discontinue the Service at any time without notice''; ``All content generated through the Service is owned by Character Technologies.'' \emph{Granted: bare licence.} \\
\addlinespace
Google (Gemini) & ``Gemini adapts to you''; ``Pick up where you left off''; cross-service personalisation. \emph{Implied: periodic tenancy.} & Right to ``modify or discontinue'' services; ``We do not make any commitments about the content within the services.'' \emph{Granted: bare licence.} \\
\bottomrule
\end{tabular}
\caption{Estate inconsistency across major AI platforms (early 2026). In each case, the marketing estate (periodic tenancy or higher) is systematically inconsistent with the ToS estate (bare licence).}
\label{tab:estate-inconsistency}
\end{table}

The pattern is consistent: marketing language --- possessive framing (``your AI''), continuity (``remembers your conversations''), and adaptation (``gets to know you'') --- signals a periodic tenancy that deepens over time and auto-renews with each session. The ToS reserves a bare licence: revocable at will, conferring no lasting interest, subject to unilateral modification without notice. The gap between these two estates is the governance failure this framework addresses.

When the marketed estate is inconsistent with the ToS estate, both English law (\emph{Street}; \emph{Autoclenz}) and American law (\emph{Bragg}; \emph{Comb}; \emph{Williams}) support disregarding the ToS characterisation in favour of the substantive relationship. The test is objective: does marketing language, product design, or system behaviour create a reasonable expectation of an estate higher than ToS grants? If yes, the substantive estate governs and Stage~1 duties attach. That conclusion establishes a tenancy, not a trust. Stage~2 still requires the additional cultivated-dependence facts identified in Section~\ref{sec:constructive-trust}. The social science literature confirms the underlying relational effect independently: \textcite{rzhang2025darkside} establish that anthropomorphic design creates pseudo-interactive spaces where transgressions in authenticity generate feelings of genuine betrayal --- estate inconsistency described in empirical rather than legal terms.

% ---------- 4.4 --------------------------------------------------------------
\subsection{California's Sliding Scale}
\label{sec:california}

California's unconscionability doctrine lowers the threshold for AI users. Under \emph{Armendariz v Foundation Health Psychcare Services, Inc.}, 24 Cal.\ 4th 83 (2000), procedural and substantive unconscionability operate on a sliding scale: neither must be present in extreme form; each compensates for the other's relative weakness. On the substantive side, the marketing-to-ToS gap represents strong unconscionability --- terms so one-sided they approach the outer limit of enforceability. On the procedural side, the architectural non-portability of personalised AI states goes beyond ordinary adhesion: the user not only had no opportunity to negotiate but has no exit that does not destroy the very estate the litigation would seek to protect. The sliding scale closes simultaneously from both ends.
